\section{Related Work}
\label{sec:related}
\emph{Side-channel model fingerprinting.} Prior work recovers neural network
architectures through microarchitectural and physical side channels, including
cache timing~\cite{TODO-cache-side-channel}, power and electromagnetic
emanations~\cite{TODO-power-em}, and GPU resource contention and memory
traces~\cite{TODO-gpu-side-channel}. These attacks require physical access to,
or co-location with, the device that executes the model, and they mostly target
inference. CHAI-FP requires only passive observation of network traffic and
targets distributed training.

\emph{Traffic analysis.} Analysis of encrypted traffic identifies websites,
applications, and user activities from packet sizes and timing~\cite{TODO-traffic-analysis},
and recent studies infer properties of machine-learning services from their
network behavior~\cite{TODO-ml-service-traffic}. These approaches classify
individual flows. CHAI-FP instead exploits the synchronized multi-client
structure of FL, in which several heterogeneous devices execute the same model
concurrently.

\emph{Attacks on federated learning.} Attacks on FL reconstruct training data
from exchanged updates~\cite{TODO-gradient-inversion}, infer membership or data
properties~\cite{TODO-fl-inference}, or poison the global
model~\cite{TODO-fl-poisoning}. They assume access to update contents, which
transport encryption and secure aggregation~\cite{TODO-secure-aggregation}
protect. CHAI-FP uses only transport metadata, which these protections leave
exposed, and its output, the identity of the model under training, can inform
the choice and configuration of such attacks.
